\chapter{Conclusion and Future Work}
\label{ch:conclusion}

\section{Summary}

This report compared Q-learning and SARSA on a warehouse navigation task in which the shortest
route runs alongside a pedestrian walkway. Both algorithms were implemented from first principles,
sharing every component except the single expression forming their update target, and were verified
against thirteen numerical checks before any experiment was run --- including an equivalence test
confirming that at $\varepsilon = 0$ they produce bit-identical action-value tables.

The comparison used thirty paired seeds with Wilcoxon signed-rank tests and bootstrap confidence
intervals, addressing directly the single-seed limitation that was the principal weakness of the
predecessor study in this repository.

\section{Findings}

\keypoint{Finding 1 --- the two algorithms separate in opposite directions on two different
questions.} Q-learning's deployed route averaged 19.40 steps against a BFS optimum of 19; SARSA's
averaged 33.34 ($p = 6.7\times10^{-7}$). Q-learning accumulated 21{,}626 walkway incidents during
training against SARSA's 5{,}068, a factor of 4.3 ($p = 1.9\times10^{-9}$). Four of five metrics
differed significantly, with large effects.

\keypoint{Finding 2 --- the mechanism is directly observable, not merely inferred.} Q-learning's
behaviour is numerically invariant to the hazard penalty: identical incident counts and identical
clearance at $-10$, $-25$, $-50$, $-100$ and $-200$. Because its target maximises over next
actions, a penalty reachable only by a non-greedy action never enters the state value. SARSA's
clearance rises monotonically over the same sweep, from 1.91 to 5.39.

\keypoint{Finding 3 --- the effect requires a localised hazard and vanishes without one.} On a
floor of scattered racking with no walkway, both algorithms converged on the identical optimum.
The distinction is a response to a property of the task, not a general ranking.

\keypoint{Finding 4 --- a four-line heuristic matched the better of the two.} The controller
recovered from the dataset achieved the optimal 19 steps with zero incidents. On this floor,
neither learned agent demonstrated superiority over it.

\section{Future work}

\keypoint{Floors where the heuristic fails.} The most valuable extension, since it is the
weakness identified in Chapter~\ref{ch:discussion}. Constructing layouts that require detours
around blocking racking would establish whether the learned policies offer anything the greedy rule
does not.

\keypoint{Expected SARSA and Double Q-learning.} Expected SARSA \cite{vanseijen2009} replaces the
sampled bootstrap with its expectation under the policy and should sit between the two algorithms
studied here, with lower variance; Double Q-learning \cite{hasselt2010} would address the
maximisation bias the $\max$ operator introduces.

\keypoint{Decaying $\varepsilon$ to zero.} Under GLIE conditions both algorithms converge to the
optimal policy. Measuring the rate at which SARSA's route shortens as $\varepsilon \to 0$ would
quantify how much of its conservatism is permanent and how much is a function of the exploration
schedule.

\keypoint{Stochastic dynamics.} Non-zero slip probability would test whether SARSA's advantage
widens when the environment, and not only the policy, is a source of randomness.

\keypoint{Function approximation.} Scaling beyond enumerable state spaces would bring the
off-policy stability question into play, which the tabular setting excludes by construction.

\section{Closing remarks}

The most useful result here is not the headline comparison, which reproduces a textbook example.
It is the sensitivity sweep, which shows Q-learning's behaviour to be entirely unmoved by a
twenty-fold change in the cost of an incident --- and the negative control, which shows that on a
task without an adjacent hazard the distinction between these two celebrated algorithms does not
exist at all.

Both point the same way. An algorithm's advantages are conditional on the task exercising the
property that algorithm addresses, and the conditions matter as much as the result.
